% GENERATED FILE. Edit canonical lessons, then run npm run generate:books.
% canonical-source-hash: fnv1a64:418c82ade8e0bc45
% canonical-lessons: ML-C371-jayil, ML-C371-tatavu, ML-C371-tatavukaran, ML-C371-vittayaykkuka, ML-C371-mocippikkuka

\chapter{Prison}
\label{ch:ml-prison}
\hlchaptermodality{371}

\begin{chapteropening}
\textbf{By the end of this chapter:} \emph{I can say a jail, imprisonment, a prisoner, to release and to free.}

Complete the last lesson of chapter 371: I can say a jail, imprisonment, a prisoner, to release and to free.
\end{chapteropening}

\section[\texorpdfstring{\ml{ജയിൽ} (jayil)}{ജയിൽ (jayil)}]{\ml{ജയിൽ} (jayil) — a jail; a prison}

\label{lesson:ML-C371-jayil}



\begin{quote}
Before the new one: say the Malayalam for to testify, then the Malayalam for to swear an oath.
\end{quote}

\subsection*{You'll want to know: \ml{ജയിൽ}}
\textbf{\ml{ജയിൽ}} — \emph{jayil} — \textquotedblleft{}a jail; a prison\textquotedblright{}.

\begin{quote}
\ml{കള്ളൻ} \ml{ഇപ്പോൾ} \ml{ജയിലിൽ} \ml{ആണ്}.
\end{quote}

\begin{quote}
\emph{kaḷḷan ippōḷ jayilil āṇŭ.}
\end{quote}

\begin{quote}
The thief is in jail now.
\end{quote}

In jail is \ml{ജയിലിൽ}.

\begin{cousinweb}
\textbf{\ml{ജയിൽ}} (\emph{jayil}) was borrowed from English \emph{jail}.
\end{cousinweb}

\begin{grammarlens}[title={What you've built}]
One more word for telling what happened.
\end{grammarlens}

\subsection*{Your turn}
\begin{itemize}
\raggedright
  \item \emph{Say it:} \emph{jayil}
  \item \emph{Say it:} \emph{jayil}, once more
  \item \emph{Say it:} \emph{kaḷḷan ippōḷ jayilil āṇŭ.} — The thief is in jail now.
  \item \emph{Recall:} say the Malayalam for to testify, then the Malayalam for to swear an oath, then say \emph{jayil} again
\end{itemize}

\begin{tcolorbox}[breakable,colback=teal!4,colframe=teal!35!black,title={Before you move on}]
What does \ml{ജയിൽ} mean? (\textquotedblleft{}A jail; a prison\textquotedblright{}.) Say it once more.
\end{tcolorbox}
\section[\texorpdfstring{\ml{തടവ്} (taṭavŭ)}{തടവ് (taṭavŭ)}]{\ml{തടവ്} (taṭavŭ) — imprisonment; custody}

\label{lesson:ML-C371-tatavu}



\begin{quote}
Before the new one: say the Malayalam for to swear an oath, then the Malayalam for a jail.
\end{quote}

\subsection*{You'll want to know: \ml{തടവ്}}
\textbf{\ml{തടവ്}} — \emph{taṭavŭ} — \textquotedblleft{}imprisonment; custody\textquotedblright{}.

\begin{quote}
\ml{കോടതി} \ml{അവന്} \ml{അഞ്ച്} \ml{വർഷം} \ml{തടവ്} \ml{വിധിച്ചു}.
\end{quote}

\begin{quote}
\emph{kōṭati avanŭ añcŭ varṣaṁ taṭavŭ vidhiccu.}
\end{quote}

\begin{quote}
The court sentenced him to five years in prison.
\end{quote}

Prison terms are counted in years: \ml{രണ്ട്} \ml{വർഷം} \ml{തടവ്}.

\begin{cousinweb}
\textbf{\ml{തടവ്}} (\emph{taṭavŭ}) is native Dravidian: DEDR 3031 lists Tamil \emph{taṭavu}, \textquotedblleft{}prison\textquotedblright{}, under a root meaning to hinder, to stop. Prison stops a person going anywhere.
\end{cousinweb}

\begin{grammarlens}[title={What you've built}]
One more word for telling what happened.
\end{grammarlens}

\subsection*{Your turn}
\begin{itemize}
\raggedright
  \item \emph{Say it:} \emph{taṭavŭ}
  \item \emph{Say it:} \emph{taṭavŭ}, once more
  \item \emph{Say it:} \emph{kōṭati avanŭ añcŭ varṣaṁ taṭavŭ vidhiccu.} — The court sentenced him to five years in prison.
  \item \emph{Recall:} say the Malayalam for to swear an oath, then the Malayalam for a jail, then say \emph{taṭavŭ} again
\end{itemize}

\begin{tcolorbox}[breakable,colback=teal!4,colframe=teal!35!black,title={Before you move on}]
What does \ml{തടവ്} mean? (\textquotedblleft{}Imprisonment; custody\textquotedblright{}.) Say it once more.
\end{tcolorbox}
\section[\texorpdfstring{\ml{തടവുകാരൻ} (taṭavukāran)}{തടവുകാരൻ (taṭavukāran)}]{\ml{തടവുകാരൻ} (taṭavukāran) — a prisoner}

\label{lesson:ML-C371-tatavukaran}



\begin{quote}
Before the new one: say the Malayalam for a jail, then the Malayalam for imprisonment.
\end{quote}

\subsection*{You'll want to know: \ml{തടവുകാരൻ}}
\textbf{\ml{തടവുകാരൻ}} — \emph{taṭavukāran} — \textquotedblleft{}a prisoner\textquotedblright{}.

\begin{quote}
\ml{തടവുകാർ} \ml{ജയിലിൽ} \ml{കൃഷി} \ml{ചെയ്യുന്നു}.
\end{quote}

\begin{quote}
\emph{taṭavukār jayilil kṛṣi ceyyunnu.}
\end{quote}

\begin{quote}
The prisoners farm in the jail.
\end{quote}

Plural: \ml{തടവുകാർ}, \emph{taṭavukār}.

\ml{തടവ്}, confinement, + -\ml{കാരൻ}, one in it.

\begin{cousinweb}
\textbf{\ml{തടവുകാരൻ}} (\emph{taṭavukāran}) is native Dravidian: \ml{തടവ്} (\emph{taṭavŭ}), \textquotedblleft{}confinement\textquotedblright{}, whose Tamil twin \emph{taṭavu}, \textquotedblleft{}prison\textquotedblright{}, stands in DEDR 3031, with the ending -\ml{കാരൻ}.
\end{cousinweb}

\begin{grammarlens}[title={What you've built}]
One more word for telling what happened.
\end{grammarlens}

\subsection*{Your turn}
\begin{itemize}
\raggedright
  \item \emph{Say it:} \emph{taṭavukāran}
  \item \emph{Say it:} \emph{taṭavukāran}, once more
  \item \emph{Say it:} \emph{taṭavukār jayilil kṛṣi ceyyunnu.} — The prisoners farm in the jail.
  \item \emph{Recall:} say the Malayalam for a jail, then the Malayalam for imprisonment, then say \emph{taṭavukāran} again
\end{itemize}

\begin{tcolorbox}[breakable,colback=teal!4,colframe=teal!35!black,title={Before you move on}]
What does \ml{തടവുകാരൻ} mean? (\textquotedblleft{}A prisoner\textquotedblright{}.) Say it once more.
\end{tcolorbox}
\section[\texorpdfstring{\ml{വിട്ടയയ്ക്കുക} (viṭṭayaykkuka)}{വിട്ടയയ്ക്കുക (viṭṭayaykkuka)}]{\ml{വിട്ടയയ്ക്കുക} (viṭṭayaykkuka) — to release; to set free}

\label{lesson:ML-C371-vittayaykkuka}



\begin{quote}
Before the new one: say the Malayalam for imprisonment, then the Malayalam for a prisoner.
\end{quote}

\subsection*{You'll want to know: \ml{വിട്ടയയ്ക്കുക}}
\textbf{\ml{വിട്ടയയ്ക്കുക}} — \emph{viṭṭayaykkuka} — \textquotedblleft{}to release; to set free\textquotedblright{}.

\begin{quote}
\ml{പോലീസ്} \ml{നിരപരാധിയെ} \ml{വിട്ടയച്ചു}.
\end{quote}

\begin{quote}
\emph{pōlīsŭ niraparādhiye viṭṭayaccu.}
\end{quote}

\begin{quote}
The police released the innocent man.
\end{quote}

Past: \ml{വിട്ടയച്ചു}, \emph{viṭṭayaccu}.

\ml{വിട്ട്}, letting go, + \ml{അയയ്ക്കുക}, to send.

\begin{cousinweb}
\textbf{\ml{വിട്ടയയ്ക്കുക}} (\emph{viṭṭayaykkuka}) is a native Dravidian compound of \ml{വിടുക} (\emph{viṭuka}), \textquotedblleft{}to let go, to release\textquotedblright{} (DEDR 5393), and \ml{അയയ്ക്കുക} (\emph{ayaykkuka}), \textquotedblleft{}to send\textquotedblright{}: to let someone go and send them off.
\end{cousinweb}

\begin{grammarlens}[title={What you've built}]
One more word for telling what happened.
\end{grammarlens}

\subsection*{Your turn}
\begin{itemize}
\raggedright
  \item \emph{Say it:} \emph{viṭṭayaykkuka}
  \item \emph{Say it:} \emph{viṭṭayaykkuka}, once more
  \item \emph{Say it:} \emph{pōlīsŭ niraparādhiye viṭṭayaccu.} — The police released the innocent man.
  \item \emph{Recall:} say the Malayalam for imprisonment, then the Malayalam for a prisoner, then say \emph{viṭṭayaykkuka} again
\end{itemize}

\begin{tcolorbox}[breakable,colback=teal!4,colframe=teal!35!black,title={Before you move on}]
What does \ml{വിട്ടയയ്ക്കുക} mean? (\textquotedblleft{}To release; to set free\textquotedblright{}.) Say it once more.
\end{tcolorbox}
\section[\texorpdfstring{\ml{മോചിപ്പിക്കുക} (mōcippikkuka)}{മോചിപ്പിക്കുക (mōcippikkuka)}]{\ml{മോചിപ്പിക്കുക} (mōcippikkuka) — to free; to liberate}

\label{lesson:ML-C371-mocippikkuka}



\begin{quote}
Before the new one: say the Malayalam for a prisoner, then the Malayalam for to release.
\end{quote}

\subsection*{You'll want to know: \ml{മോചിപ്പിക്കുക}}
\textbf{\ml{മോചിപ്പിക്കുക}} — \emph{mōcippikkuka} — \textquotedblleft{}to free; to liberate\textquotedblright{}.

\begin{quote}
\ml{പോലീസ്} \ml{കുട്ടിയെ} \ml{മോചിപ്പിച്ചു}.
\end{quote}

\begin{quote}
\emph{pōlīsŭ kuṭṭiye mōcippiccu.}
\end{quote}

\begin{quote}
The police freed the child.
\end{quote}

Past: \ml{മോചിപ്പിച്ചു}, \emph{mōcippiccu}.

The noun is \ml{മോചനം}, release, freedom.

\begin{cousinweb}
\textbf{\ml{മോചിപ്പിക്കുക}} (\emph{mōcippikkuka}) is built from Sanskrit \emph{mocana}, \textquotedblleft{}releasing, setting free\textquotedblright{} (Monier-Williams), with the Malayalam causative ending -\ml{ിപ്പിക്കുക}.
\end{cousinweb}

\begin{grammarlens}[title={What you've built}]
One more word for telling what happened.
\end{grammarlens}

\subsection*{Your turn}
\begin{itemize}
\raggedright
  \item \emph{Say it:} \emph{mōcippikkuka}
  \item \emph{Say it:} \emph{mōcippikkuka}, once more
  \item \emph{Say it:} \emph{pōlīsŭ kuṭṭiye mōcippiccu.} — The police freed the child.
  \item \emph{Recall:} say the Malayalam for a prisoner, then the Malayalam for to release, then say \emph{mōcippikkuka} again
\end{itemize}

\begin{tcolorbox}[breakable,colback=teal!4,colframe=teal!35!black,title={Before you move on}]
What does \ml{മോചിപ്പിക്കുക} mean? (\textquotedblleft{}To free; to liberate\textquotedblright{}.) Say it once more.
\end{tcolorbox}
